% SEGMENT OLP-0128-S01
% Part: first-order-logic
% Chapter: completeness
% Section: outline

\documentclass[../../../include/open-logic-section]{subfiles}

\begin{document}

\iftag{FOL}
      {\olfileid{fol}{com}{out}}
      {\olfileid{pl}{com}{out}}

\olsection{நிறுவலின் உருவரை}

% SEGMENT OLP-0128-S02
நிறைவு தேற்றத்தின் நிறுவல் சற்றுச் சிக்கலானது; அதை முதன்முறையாகப்
படிக்கும்போது வழி தவறுவது எளிது. ஆகவே, நிறுவலின் உருவரையை முதலில்
காண்போம். நிறுவலுக்கான வழியைக் காண உதவும் ஒரு பார்வை மாற்றமே முதல் படி.
நிறைவை, ``$\Gamma \Entails !A$ எனில் $\Gamma \Proves !A$'' என்று
கருதும்போது, ஓர் எண்ணத்தை உருவாக்குவதுகூடக் கடினமாக இருக்கலாம்:
$\Gamma \Proves !A$ என்பதைக் காட்ட நாம் !!a{derivation} ஒன்றைக் கண்டுபிடிக்க
வேண்டும்; ஆனால் $\Gamma \Entails !A$ என்ற கருதுகோள் அதற்கு எந்தவகையிலும்
உதவுவது போலத் தெரியவில்லை. சில நிறுவல் அமைப்புகளில் !!a{derivation} ஐ
நேரடியாகக் கட்டமைக்க முடியும்; நாம் சற்றே வேறொரு அணுகுமுறையைப் பின்பற்றுவோம்.
தேவையான பார்வை மாற்றம் இதுதான்: ``$\Gamma$ முரணற்றது எனில், அது
நிறைவுறத்தக்கது'' என்றும் நிறைவை வகுக்கலாம். $\Gamma$ விலுள்ள தகவலையும் அது
முரணற்றது என்ற கருதுகோளையும் சேர்த்துப் பயன்படுத்தி, $\Gamma$ விலுள்ள ஒவ்வொரு
\iftag{FOL}{!!{sentence}}{!!{formula}} ஐயும் நிறைவுறுத்தும்
\iftag{FOL}{!!a{structure}}{!!a{valuation}} ஐ நாம் கட்டமைக்க முடியலாம்.
எப்படிப்பட்ட \iftag{FOL}{!!{structure}}{!!{valuation}} தேவை என்பது நமக்குத்
தெரியும்: $\Gamma$ விவரிப்பதைப் போலவே உள்ள ஒன்று!{}

% SEGMENT OLP-0128-S03
$\Gamma$ வில் \iftag{FOL}{அணு !!{sentence}s}{!!{propositional variable}s}
மட்டுமே இருந்தால், அதற்கான மாதிரியைக் கட்டமைப்பது எளிது.\iftag{FOL}{ அந்த அணு
  !!{sentence}s அனைத்தும் $\Atom{P}{a_1,\dots,a_n}$ என்ற வடிவில் உள்ளன என்றும்,
  இதில் $a_i$ கள் !!{constant}s என்றும் கொள்க.}{} நாம் செய்ய வேண்டியது
\iftag{FOL}{%
  !!a{domain}~$\Domain{M}$ ஐயும், $\Sat{M}{\Atom{P}{a_1,\ldots,a_n}}$
  ஆகுமாறு $P$ க்கான ஓர் ஒதுக்கீட்டையும் தருவதுதான். இது கடினமல்ல:
  $\Domain{M} = \Nat$, $\Assign{\Obj c_i}{M} = i$ என வைக்க; மேலும் ஒவ்வொரு
  $\Atom{P}{a_1,\ldots,a_n} \in \Gamma$ க்கும், $\tuple{k_1, \dots, k_n}$
  ஐ $\Assign{P}{M}$ இல் வைக்க; இங்கு $k_i$ என்பது $a_i$ என்ற மாறிலிக்
  குறியின் சுட்டெண் (அதாவது, $a_i \ident \Obj c_{k_i}$).}{%
  எல்லா $p \in \Gamma$ க்கும் $\pSat{v}{p}$ ஆகுமாறு !!a{valuation}~$\pAssign{v}$
  ஐத் தருவதுதான். $\pAssign{v}(p) = \True$ அப்போதும் அப்போது மட்டுமே
  $p \in \Gamma$ என வைப்போம்.}

% SEGMENT OLP-0128-S04
இப்போது $\Gamma$ வில் $\lnot !B$ என்ற ஏதோ !!{formula} இருப்பதாகவும், $!B$
அணுவாய்பாடு என்றும் கொள்க. $\iftag{FOL}{\Struct{M}}{\pAssign{v}}$ இன்
கட்டமைப்பு $\lnot !B$ ஐ மெய்யாக்க இயலாமல் செய்துவிடுமோ என்று நாம் கவலைப்படலாம்.
இங்குதான்~$\Gamma$ வின் முரணின்மை உதவுகிறது: $\lnot !B \in \Gamma$ எனில்,
$!B \notin \Gamma$; இல்லையெனில் $\Gamma$ முரணுடையதாகிவிடும். மேலும்
$!B \notin \Gamma$ எனில், நமது
$\iftag{FOL}{\Struct{M}}{\pAssign{v}}$ கட்டமைப்பின்படி
$\iftag{FOL}{\Sat/{M}}{\pSat/{v}}{!B}$; ஆகவே
$\iftag{FOL}{\Sat{M}}{\pSat{v}}{\lnot !B}$. இதுவரை எல்லாம் சரியாக உள்ளது.

% SEGMENT OLP-0128-S05
$\Gamma$ வில் சிக்கலான, அணுவல்லாத வாய்பாடுகள் இருந்தால் என்ன செய்வது?
$!A \land !B$ அதில் உள்ளது என்க. அதை மெய்யாக்க, $!A$, $!B$ இரண்டும்
$\Gamma$ வில் இருப்பதுபோல் செயல்பட வேண்டும். $!A \lor !B \in \Gamma$ எனில்,
அவற்றில் குறைந்தது ஒன்றை மெய்யாக்க வேண்டும்; அதாவது, அவற்றில் ஒன்று
$\Gamma$ வில் இருப்பதுபோல் செயல்பட வேண்டும்.

% SEGMENT OLP-0128-S06
இது பின்வரும் எண்ணத்தைத் தருகிறது: (a)~கிடைக்கும் கணத்தை முரணற்றதாக வைத்திருக்கவும்,
(b)~சாத்தியமான ஒவ்வோர் அணு !!{sentence}~$!A$ க்கும், $!A$ அல்லது $\lnot !A$
கிடைக்கும் கணத்தில் இருப்பதை உறுதிசெய்யவும், (c)~$!A \land !B$ கணத்தில்
இருக்கும்போதெல்லாம் $!A$, $!B$ இரண்டும் அதில் இருக்கவும், $!A \lor !B$
கணத்தில் இருந்தால் $!A$, $!B$ இவற்றில் குறைந்தது ஒன்றாவது அதில் இருக்கவும்,
இதுபோலத் தொடரவும் ஏற்ற கூடுதல் !!{formula}s ஐ~$\Gamma$ வுடன் சேர்க்கிறோம்.
இதை (முடிவில்லாமல்கூட) தொடர்ந்து செய்கிறோம். இவ்வாறு சேர்க்கப்பட்ட எல்லா
!!{formula}s இன் கணத்தை~$\Gamma^*$ என அழைப்போம். மேலுள்ள நமது கட்டமைப்பு,
\iftag{FOL}{!!a{structure}~$\Struct{M}$}{!!a{valuation}~$\pAssign{v}$} ஐத்
தரும். அது~$\Gamma^*$ விலுள்ள எல்லா வாக்கியங்களையும் நிறைவுறுத்துகிறது
என்பதைத் தொகுத்தறிதலால் நிறுவ முடியும்; $\Gamma \subseteq \Gamma^*$ ஆதலால்,
அது~$\Gamma$ விலுள்ள எல்லா வாக்கியங்களையும் நிறைவுறுத்துகிறது. (a), (b)
ஆகியவற்றை உறுதிசெய்வதே போதும் எனத் தெரியவருகிறது. (b) ஐ நிறைவுசெய்யும்
வாக்கியக் கணத்தை \emph{நிறைவானது} என்கிறோம். ஆகவே, முரணற்ற
கணம்~$\Gamma$ ஐ முரணற்ற, நிறைவான கணம்~$\Gamma^*$ ஆக விரிவாக்குவதே நமது பணி.

% SEGMENT OLP-0128-S07
\iftag{FOL}{%
இந்தத் திட்டத்தில் ஒரு சிக்கல் உள்ளது: $\lexists[x][!A(x)] \in \Gamma$ எனில்,
ஏதோ !!{constant}~$c$ ஐத் தேர்ந்து, இந்தச் செயல்முறையில் $!A(c)$ ஐச் சேர்க்க
விரும்புவோம். இதை எப்போதும் செய்ய முடியும் என்பது எப்படித் தெரியும்? நமது
மொழியில் சில !!{constant}s மட்டுமே இருந்து, அவை ஒவ்வொன்றுக்கும்
$\lnot !A(c) \in \Gamma$ ஆக இருக்கலாம். $!A(c)$ ஐயும் சேர்க்க முடியாது;
அப்படிச் செய்தால் கணம் முரணுடையதாகிவிடும்; $\Struct{M}$ ஆனது $!A(c)$ ஐயா,
$\lnot !A(c)$ ஐயா மெய்யாக்க வேண்டும் எனத் தீர்மானிக்க இயலாது. மேலும்,
$\Gamma$ ஆனது மாறிலிக் குறிகளே இல்லாத ஒரு மொழியின் வாக்கியங்களை மட்டுமே
கொண்டிருக்கலாம் (எடுத்துக்காட்டாக, கணக்கோட்பாட்டின் மொழி).

% SEGMENT OLP-0128-S08
இந்தச் சிக்கலுக்கான தீர்வு எளியது: தொடக்கத்திலேயே முடிவிலா எண்ணிக்கையிலான
மாறிலிக் குறிகளையும், அவற்றை அளவையடைகளுடன் சரியான முறையில் இணைக்கும்
வாக்கியங்களையும் சேர்த்துவிடுகிறோம். (இது ஒரு முரண்பாட்டை அறிமுகப்படுத்தாது
என்பதை நிச்சயமாக நாம் சரிபார்க்க வேண்டும்.)

% SEGMENT OLP-0128-S09
அணு வாக்கியங்களில் !!{constant}s மட்டுமே இருந்தால் நமது முதற் கட்டமைப்பு
நன்றாகச் செயல்படுகிறது. ஆனால் மொழியில் !!{function}s உம் இருக்கலாம். அப்போது,
எல்லாம் செயல்படுமாறு அந்த !!{function}s க்கு ஒதுக்க வேண்டிய சரியான
$\Nat$-மீதான சார்புகளைக் கண்டுபிடிப்பது கடினமாக இருக்கலாம். எனவே இன்னொரு
உத்தியைப் பயன்படுத்துவோம்: $\Obj c_i$ ஐ விளக்க $i$ ஐப் பயன்படுத்துவதற்குப்
பதிலாக, !!{constant}s இன் கணத்தையே பொருட்களமாக எடுத்துக்கொள்வோம். அப்போது
$\Struct M$ ஒவ்வொரு !!{constant} ஐயும் அதற்கே ஒதுக்கலாம்:
$\Assign{\Obj c_i}{M} = \Obj c_i$. இதை இன்னும் முழுமையாகச் செய்யலாம்:
$\Domain{M}$ ஐ மொழியின் எல்லாச் \emph{சொற்களின்} கணமாகவே கொள்வோம்!{} இவ்வாறு
செய்தால், சொற்களை உள்ளீடுகளாகவும் சொற்களை மதிப்புகளாகவும் கொண்ட சார்புகளை
!!{function}s க்கு ஒதுக்க ஒரு வெளிப்படையான முறை கிடைக்கிறது: $n$ சொற்கள்
$t_1$, \dots,~$t_n$ உள்ளீடாகத் தரப்படும்போது $\Obj f^n_i(t_1, \dots, t_n)$
என்ற சொல்லை மதிப்பாகத் தரும் சார்பை !!{function}~$\Obj f^n_i$ க்கு
ஒதுக்குகிறோம்.

% SEGMENT OLP-0128-S10
புதிரின் கடைசிப் பகுதி~$\eq$ ஐ எவ்வாறு கையாள்வது என்பதாகும்.
!!{predicate}~$\eq$ க்கு நிலையான பொருள்விளக்கம் உண்டு:
$\Sat{M}{\eq[t][t']}$ அப்போதும் அப்போது மட்டுமே
$\Value{t}{M} = \Value{t'}{M}$. ஒரு சொல்~$t$ இன் !!{value} $t$ தானே ஆகுமாறு
அமைத்தால், $t$, $t'$ ஆகியவை ஒரே சொல்லாக இல்லாதவரை $\eq[t][t']$ வடிவிலுள்ள
\emph{எந்த} வாக்கியத்தையும் இந்த !!{structure} மெய்யாக்காது. இது ஒரு சிக்கலே;
ஏனெனில் !!{function}s கொண்ட மொழியில் அமைந்த பெரும்பாலும் ஒவ்வொரு சுவையான
கோட்பாட்டிலும், ஒரே சொல்லாக இல்லாத $t$, $t'$ களுக்கான $\eq[t][t']$ என்ற
தேற்றங்கள் இருக்கும் (எடுத்துக்காட்டாக, எண்கணிதக் கோட்பாடுகளில்:
$\eq[(\Obj 0+ \Obj 0)][\Obj 0]$). இதைத் தீர்க்க~$\Struct M$ இன் பொருட்களத்தை
மாற்றுகிறோம்: $\Domain{M}$ இல் சொற்களைப் பொருள்களாகப் பயன்படுத்துவதற்குப்
பதிலாக, சொற்களின் கணங்களைப் பயன்படுத்துகிறோம்; ஒவ்வொரு கணமும்~$\Gamma$
விலுள்ள வாக்கியங்கள் ஒன்றானவை எனக் கோரும் எல்லாச் சொற்களையும் கொண்டிருக்கும்.
எடுத்துக்காட்டாக, $\Gamma$ ஓர் எண்கணிதக் கோட்பாடு எனில், இந்தக் கணங்களில்
ஒன்று $\Obj 0$, $(\Obj 0 + \Obj 0)$, $(\Obj 0 \times \Obj 0)$ முதலியவற்றைக்
கொண்டிருக்கும். இந்தக் கணத்தையே $\Obj 0$ க்கு ஒதுக்குவோம்; அந்தக் கணத்திலுள்ள
எல்லாச் சொற்களுக்கும், எடுத்துக்காட்டாக $(\Obj 0 + \Obj 0)$ க்கும், இதுவே
மதிப்பாக அமையும். எனவே, திருத்திய இந்த !!{structure} இல்
$\eq[(\Obj 0+ \Obj 0)][\Obj 0]$ என்ற வாக்கியம் மெய்யாகும்.}{}

% SEGMENT OLP-0128-S11
இப்போது நாம் செய்யப்போவதைத் தொகுத்துக் கூறலாம். முதலில் !!{complete} முரணற்ற
கணங்களின் பண்புகளை ஆராய்வோம். குறிப்பாக, !!a{complete} முரணற்ற கணம்
$!A \land !B$ ஐக் கொண்டிருக்கும் அப்போதும் அப்போது மட்டுமே அது $!A$,~$!B$
இரண்டையும் கொண்டிருக்கும்; $!A \lor !B$ ஐக் கொண்டிருக்கும் அப்போதும் அப்போது
மட்டுமே அவற்றில் குறைந்தது ஒன்றைக் கொண்டிருக்கும்; இதுபோல மற்ற
இணைப்புகளுக்கும் அமையும் என்பதை நிறுவுவோம் (\olref[ccs]{prop:ccs}).\iftag{FOL}{
  அடுத்து, ``செறிவுற்ற'' வாக்கியக் கணங்களை வரையறுத்து ஆராய்வோம். ஒவ்வொரு
  அளவையிட்ட !!{sentence} ஐயும் அதன் நிகழ்வுகளுடன் இணைக்கும் நிபந்தனைக்
  கூற்றுகளைக் கொண்ட கணமே செறிவுற்ற கணம் (\olref[hen]{defn:henkin-exp}).
  எந்த முரணற்ற கணம்~$\Gamma$ ஐயும் ஒரு செறிவுற்ற கணம்~$\Gamma'$ ஆக
  எப்போதும் விரிவாக்க முடியும் என்பதைக் காட்டுவோம் (\olref[hen]{lem:henkin}).
  ஒரு கணம் முரணற்றதாகவும், செறிவுற்றதாகவும், !!{complete} ஆகவும் இருந்தால்,
  அது \iftag{prvEx}{$\lexists[x][!A(x)]$ ஐக் கொண்டிருக்கும் அப்போதும் அப்போது
    மட்டுமே ஏதோ மூடிய சொல்~$t$ க்கு $!A(t)$ ஐக் கொண்டிருக்கும்}{}\iftag{defEx,defAll}{}{
  மேலும் }\iftag{prvAll}{$\lforall[x][!A(x)]$ ஐக் கொண்டிருக்கும் அப்போதும்
    அப்போது மட்டுமே எல்லா மூடிய சொற்கள்~$t$ க்கும் $!A(t)$ ஐக் கொண்டிருக்கும்}{}
  என்ற பண்பையும் பெறும் (\olref[hen]{prop:saturated-instances}).}{}
பிறகு, \iftag{FOL}{செறிவுற்ற}{} முரணற்ற கணம்~$\iftag{FOL}{\Gamma'}{\Gamma}$
ஐ ஒரு \iftag{FOL}{செறிவுற்ற, முரணற்ற,}{முரணற்ற} !!{complete}
கணம்~$\Gamma^*$ ஆக விரிவாக்கலாம் என்பதைக் காட்டுவோம்
(\olref[lin]{lem:lindenbaum}). இந்த~$\Gamma^*$ ஐப் பயன்படுத்தி நமது
\iftag{FOL}{சொல் மாதிரி~$\Struct M(\Gamma^*)$}{!!{valuation}~$\pAssign v(\Gamma^*)$}
ஐ வரையறுப்போம். \iftag{FOL}{சொல் மாதிரியின் பொருட்களம் மூடிய சொற்களின்
  கணமாகும்; அதன் !!{predicate}s இன் பொருள்விளக்கத்தைக் கொடுக்கும் அணு
  !!{sentence}s}{மெய்மதிப்பு ஒதுக்கீட்டைக் கொடுக்கும் !!{propositional
    variable}s}~$\Gamma^*$ வில் உள்ளன (\olref[mod]{defn:termmodel}).
\iftag{FOL}{செறிவுற்ற,}{} நிறைவான முரணற்ற கணங்களின் பண்புகளைப் பயன்படுத்தி,
$\iftag{FOL}{\Sat{M(\Gamma^*)}}{\pSat{v(\Gamma^*)}}{!A}$ அப்போதும் அப்போது
மட்டுமே $!A \in \Gamma^*$ என்பதை நிறுவுவோம் (\olref[mod]{lem:truth});
ஆகவே குறிப்பாக,
$\iftag{FOL}{\Sat{M(\Gamma^*)}}{\pSat{v(\Gamma^*)}}{\Gamma}$.
\iftag{FOL}{இறுதியாக, $\Gamma$ ஆனது~$\eq$ ஐயும் கொண்டிருக்கும்போது சொல்
  மாதிரியை எவ்வாறு வரையறுப்பது என்று ஆராய்ந்து
  (\olref[ide]{defn:term-model-factor}), அது~$\Gamma^*$ ஐ நிறைவுறுத்துகிறது
  என்பதைக் காட்டுவோம் (\olref[ide]{lem:truth}).}{}

\end{document}
